\chapter{The separably closed rank-one cover}
\label{route-ch-sepclosed}

Let \(k\) now be separably closed and retain H1--H2 for \(C/k\).
Separably closedness is used to construct translators and prove coverage; it
is not inserted into the rank-one Abel isomorphism itself.

\section{Translated charts}

\begin{definition}[The exact translator and chart family]
  \label{route-def-translated-charts}
  \lean{AlgebraicGeometry.PicRankOneTranslatorIndex,
    AlgebraicGeometry.pic0GenusTranslationIso,
    AlgebraicGeometry.picRankOneTranslatedChart,
    AlgebraicGeometry.picRankOneTranslatedChart_isOpenImmersion}
  \leanok
  \dcref{custom}
  \uses{route-thm-rank-one-abel-iso}

  Put \(\mathcal P^0_\Sigma=(\operatorname{pic0SigmaSheaf} C).1\).  A
  translator is exactly an element of the Lean subtype
  \[
    \operatorname{PicRankOneTranslatorIndex}(C)
      =\left\{a:(C\mathbin\otimes\operatorname{overSpec}(k,k)).\mathrm{left}.
          \mathrm{CechPic}\ \middle|\
          \operatorname{classDeg}_k(a)=g\right\}.
  \]
  On the slice site, multiplication by its theta family gives the natural
  isomorphism
  \[
    \tau_a(T):\operatorname{Pic}^0_C(T)
       \xrightarrow{\sim}\operatorname{Pic}^{g}_C(T),
    \qquad \lambda\longmapsto\lambda\otimes a_T .
  \]
  With \(W\) the represented rank-one divisor open from
  Theorem~\ref{route-thm-rank-one-open}, the actual big-site chart is
  \[
    j_a:\operatorname{yoneda}(W.\mathrm{left})
      \longrightarrow\mathcal P^0_\Sigma .
  \]
  It is the canonical rank-one Abel isomorphism, followed by the rank-one-locus
  inclusion, followed by the inverse of the sigma extension of \(\tau_a\).
  Every \(j_a\) is relatively represented by open immersions.  The later
  representability theorem uses the canonical sigma-extension-to-slice bridge
  to recover the target \(\operatorname{pic0TypeFunctor}(C)\); the chart target
  itself is not definitionally that slice functor.
\end{definition}

\begin{proof}
  The three factors have the exact Lean carriers
  \[
    \operatorname{yoneda}(W.\mathrm{left})\xrightarrow{\sim}
    \operatorname{rankOneLocus}(\pi)
      \longrightarrow
    \operatorname{sigmaExtension}(\operatorname{Pic}^{g}_C)
      \xrightarrow{\operatorname{sigmaExtension}(\tau_a)^{-1}}
    \mathcal P^0_\Sigma .
  \]
  The first and last factors are natural isomorphisms and the middle factor is
  the checked relatively
  represented open inclusion.  Closure of open immersions under composition
  proves the claim.
\end{proof}

\begin{theorem}[Translated rank-one coverage]
  \label{route-thm-translated-coverage}
  \lean{AlgebraicGeometry.exists_sepClosedTranslated_mem_picRankOneOpen,
    AlgebraicGeometry.exists_picRankOneTranslatedChart_fieldFactorization,
    AlgebraicGeometry.picRankOneTranslatedChart_pointwiseCoverage}
  \leanok
  \source{abelian-varieties:page-0107, stacks-project}
  \dcref{Milne III, Lemma 5.2(b); Stacks Tag 056U}
  \uses{route-def-translated-charts}

  The auxiliary producer has the following exact field-valued type.  For every
  field extension \(K/k\) and every
  \(\mu\in\operatorname{pic0Subgroup}(C,\operatorname{overSpec}(k,K))\), it
  returns a
  finite separable extension \(L/K\), the scalar-tower instances
  \(k\to K\to L\), a finite-support datum on \(C_L\), its
  residue-degree-one certificate, and the Cech subtraction equality; its final
  conclusion is that the translated layer element belongs to the rank-one open.
  In symbols, its witness has the shape
  \[
    \begin{aligned}
    \exists L\;&[\operatorname{Field}L]\,[\operatorname{Algebra}kL]
      [\operatorname{Algebra}KL]\,[\operatorname{IsScalarTower}kKL]\\
      & [\operatorname{Module.Finite}_K L]
      [\operatorname{Algebra.IsSeparable}KL]\;\exists d\;r,\\
      &\operatorname{classDeg}_k(\operatorname{chartTwistClass}(d.m,r.Z))=g
        \;\wedge\;\operatorname{Finite}(r.S.\operatorname{support})\\
      &\wedge\;\bigl(\forall p\in r.S.\operatorname{support},
        \operatorname{residueDeg}_L(p)=1\bigr)\\
      &\wedge\;\operatorname{picClass}_L(r.S)=
        \operatorname{CechPic.map}(C,k,L)
          (\operatorname{chartTwistClass}(0,r.Z)^{-1})\\
      &\wedge\;\operatorname{rankOneLayer}(\mu,d)\in
        \operatorname{PicRankOneOpen}(\pi)(\operatorname{overSpec}(k,K)).
    \end{aligned}
  \]
  The public theorem
  \texttt{exists\_picRankOneTranslatedChart\_fieldFactorization} deliberately
  hides that auxiliary \(L,d,r\): its exact conclusion is only
  \[
    \begin{aligned}
      &\exists a:\operatorname{PicRankOneTranslatorIndex}(C)\;\exists x:
        (\operatorname{overSpec}(k,K)).\mathrm{left}
          \longrightarrow W.\mathrm{left},\\
      &(j_a)_{(\operatorname{overSpec}(k,K)).\mathrm{left}}(x)
        =\langle(\operatorname{overSpec}(k,K)).\mathrm{hom},\mu\rangle\\
      &\hspace{8em}\in
        \mathcal P^0_\Sigma((\operatorname{overSpec}(k,K)).\mathrm{left}).
    \end{aligned}
  \]
  Thus \(L\) is not a field-valued output of the public factorization and is
  never identified with \(k\) or \(K\).  Consequently, for every \(k\)-scheme
  \(T\), every \(s\in\mathcal P^0_\Sigma(T)\), and every point \(t\in T\), there are
  an open neighbourhood \(U\subseteq T\), an index \(a\), and
  \(x_U:U\to W.\mathrm{left}\) with \(j_a(x_U)=s|_U\).  Equivalently, the family
  \(\{j_a\}_a\) is pointwise covering and its sigma-desc map is
  Zariski-locally surjective.
\end{theorem}

\begin{proof}
  Restrict the slice component of \(s\) to \(\kappa(t)\).  The translated-drop construction first
  produces the displayed finite separable extension \(L/\kappa(t)\), a finite
  support on \(C_L\) of residue degree one over \(L\), and the Cech
  subtraction equality.  The resulting translated class is then proved to
  lie in the public \(\operatorname{PicRankOneOpen}\), and the canonical Abel
  inverse gives the factorization through \(W\).  There is no step in which
  separable closedness identifies the auxiliary residue field with \(k\), and
  no descent of the witness from \(L\) (or from \(\kappa(t)\)) to \(k\) is
  claimed.  Relative representability of \(j_a\) pulls this residue-field factorization
  back to an open neighbourhood of \(t\).  Stacks Tag 056U is used only for
  density of closed points with finite separable residue field on a smooth
  scheme; obtaining the required rational point over a separably closed field
  is a separate specialization.

  Milne's Lemma 5.2(b) is the classical fibrewise model: a general
  degree-\(g\) divisor has \(h^0=1\).  The Lean proof is stronger in the
  direction needed here because it constructs the translator and proves
  membership in the exact etale-local rank-one predicate.  Milne's surrounding
  Section III.5 assumes a complete nonsingular curve of positive genus and a
  chosen rational point; the citation is only the fibrewise model, while the
  checked translated-drop theorem retains AJCR's displayed hypotheses and
  introduces no such point.
\end{proof}

\section{The exact separably closed representer}

\begin{theorem}[Representability from the translated atlas]
  \label{route-thm-sepclosed-representability}
  \lean{AlgebraicGeometry.pic0_sepClosed_representableBy_of_isOpen,
    AlgebraicGeometry.pic0_sepClosed_representableBy}
  \leanok
  \source{stacks-project, kleiman-picard}
  \dcref{Tag 01JJ; Theorem 4.8(1)}
  \uses{route-thm-translated-coverage}

  There is one selected pair
  \[
    (J_k,r_k),\qquad J_k\longrightarrow\Spec k,\qquad
    r_k:h_{J_k}\xrightarrow{\sim}\operatorname{Pic}^0_C,
  \]
  where \(J_k\) is obtained by gluing the represented open charts
  \(\{W\xrightarrow{j_a}\operatorname{Pic}^0_C\}_a\).
  In Lean this pair is the dependent value
  \[
    \operatorname{pic0\_sepClosed\_representableBy}(C):
    \sum_{J:\Over(\Spec k)}
      \operatorname{RepresentableBy}(\operatorname{pic0TypeFunctor}(C),J).
  \]
\end{theorem}

\begin{proof}
  Stacks Tag 01JJ says that a Zariski sheaf covered by representable
  subfunctors whose inclusions are relatively represented by open immersions
  is representable.  Apply it to the charts of
  Definition~\ref{route-def-translated-charts}: openness is proved there and
  coverage is Theorem~\ref{route-thm-translated-coverage}.  The overlap maps
  are forced by equality of their universal Picard classes, so the glued
  functor of points is naturally \(\operatorname{Pic}^0_C\).
  Kleiman's Theorem 4.8 supplies the surrounding classical Picard
  representability theorem, but the checked Lean result is specifically the
  translated rank-one atlas specialization just described.
\end{proof}

\begin{theorem}[Finiteness and same-carrier data over a separably closed field]
  \label{route-thm-sepclosed-finiteness}
  \lean{AlgebraicGeometry.locallyOfFiniteType_pic0_sepClosed_representableBy,
    AlgebraicGeometry.locallyOfFinitePresentation_pic0_sepClosed_representableBy,
    AlgebraicGeometry.quasiCompact_pic0SepClosedRepresenter,
    AlgebraicGeometry.quasiSeparatedSpace_pic0SepClosedRepresenter,
    AlgebraicGeometry.picRepDatumSepClosed,
    AlgebraicGeometry.jacobianDataSepClosed}
  \leanok
  \source{kleiman-picard, stacks-project}
  \dcref{Kleiman Theorem 4.8(1); Stacks Tag 01TX}
  \uses{route-thm-sepclosed-representability, route-rem-admissible-source}

  The selected structure morphism \(J_k\to\Spec k\) is locally of finite
  presentation and quasi-compact, and \(J_k\) is quasi-separated.  The exact
  same pair \((J_k,r_k)\) defines
  \[
    \operatorname{PicRepDatum}(k,k,C)
    \quad\text{and}\quad
    \operatorname{JacobianData}(C);
  \]
  the second adds the quasi-compactness certificate, while neither handoff
  changes the carrier, representation, group object, or universal element.
\end{theorem}

\begin{proof}
  Each rank-one chart is locally of finite type, so the glued structure map
  is locally of finite type.  Since \(\Spec k\) is Noetherian, Stacks
  Tag 01TX upgrades this to local finite presentation.  The represented
  admissible divisor source is quasi-compact and its
  Abel morphism to this exact \(J_k\) is etale-locally surjective, hence
  surjective on underlying points; quasi-compactness descends through that
  map.  Local finite type over a field makes \(J_k\) locally Noetherian and
  therefore quasi-separated in the checked interface.  Finally the
  definitional projection lemmas for
  \(\operatorname{picRepDatumSepClosed}\) and
  \(\operatorname{jacobianDataSepClosed}\) identify both fields with
  \((J_k,r_k)\) literally.
\end{proof}
